\section{Introduction}
\label{sec:introduction}

Many canonical discussions of physical computation and reversible logic begin from a prior vocabulary of bits, registers, and truth tables \cite{landauer1961,bennett1973,fredkintoffoli1982conservative,bennett1982thermodynamics}. That question is important, but it already assumes the existence of a symbolic layer. Before implementation lies a harder prior question: when does a coarse-grained layer deserve to be described in logical terms at all? A system may admit labels that look like \(0\) and \(1\) without supporting stable predicates, reliable negation, or closed gate operators. In that case, logic has been projected onto the system rather than discovered in it.

This paper gives an operational answer. We work in the Six Birds emergence calculus \cite{tsiokos2026sixbirds}, whose six primitives are P1 operator rewrite, P2 gating / constraints, P3 autonomous protocol holonomy, P4 sectors / invariants, P5 packaging, and P6 accounting / audit. Our central claim is that logic is not a seventh primitive and not a hidden Platonic layer. It is a layer-relative achievement of these six ingredients. A system has a logic layer when packaged binary predicates become stable, when induced macro-operators close on those predicates, and when the reliability and loss of that closure can be audited. In one sentence: a logic layer is an audited closure of packaged predicates under induced dynamics.

That perspective shifts the usual starting point. In the feasibility view, propositions are not abstract atoms but definable feasible sets, and logical connectives are stable closure operations on those sets. In the operator view, bits are not fundamental carriers but packaged equivalence classes of microstates, and gates are macro-updates induced by the underlying dynamics. The two views are complementary. The first gives semantics to truth; the second tells us when those semantics are executable by a substrate. Together they replace the familiar question ``how is logic implemented?'' with the more basic question ``when has logic emerged strongly enough to close?'' \cite{pawlak1982rough,ellerman2010partitions,shalizi2025macrostate}

The title emphasizes XOR for a reason. Parity-like variables aligned with a system's invariant or sector structure are often especially natural coarse variables; in this paper parity is the controlled test case for that more general idea \cite{simonando1961aggregation,shalizi2025macrostate,barnettseth2023dynamical}. It survives cancellation of micro-detail, aligns naturally with invariant structure, and persists under coarse-graining in regimes where more brittle variables do not. That does not make XOR metaphysically privileged. It makes parity an empirical test case for the broader thesis of the paper: some logical forms are easier for substrates to stabilize than others. The right question is therefore not whether a substrate can be made to emulate a desired truth table in principle, but which logical distinctions it supports natively, robustly, and at what accounting cost.

Our emphasis is therefore diagnostic rather than merely implementational. We introduce a suite of quantitative tests for candidate logic layers: predicate stability, route mismatch, closure defect, gate error, conditional entropy, unretained input information, and entropy-production audits. These tests distinguish genuine closure from symbolic overinterpretation. They let us ask whether two routes to the same purported macro-description agree, whether a packaged bit persists long enough to matter, whether a gate is really a closed macro-operator rather than a protocol artifact, and whether information that appears to disappear has in fact been erased, hidden, or retained in side variables. This is also where the Six Birds grammar becomes empirically useful: packaging makes predicates possible, constraints give them meaning, protocol effects create route dependence, sectors stabilize them, rewrite repairs or redesigns closure, and accounting determines what can be maintained without cheating.

Our starting point is therefore closer to coarse-graining and audited closure than to implementation-first accounts of logic. Internally, the paper draws on the Six Birds foundations \cite{tsiokos2026sixbirds} and on the broader closure-and-substrate program developed in \cite{tsiokos2026countstone,tsiokos2026minsubstrate}. Externally, the descent of dynamics to coarse variables is related to lumpability and effective Markov aggregation \cite{kemeny1976finite,buchholz1994lumpability,simonando1961aggregation,givon2004extracting}, while the distinction between retained and erased descriptions connects directly to the classic literature on irreversibility and reversible computation \cite{landauer1961,bennett1973,bennett1982thermodynamics}.

The payoff is a controlled program rather than a metaphor. We build finite Markov laboratories in which NOT, AND, and reversible XOR embeddings can be realized with tunable noise, metastability, and coarse-graining. Within those laboratories we ask four concrete questions. First, which binary coarse variables close most robustly under the induced dynamics? Second, how do familiar gates behave as noise, staging barriers, and timescale stress are varied? Third, what changes when a reversible embedding is coarse-grained to an apparently simpler but information-losing gate? Fourth, can stable bits and gates be recovered from transition structure without declaring them by hand in advance?

This paper answers all four. Parity beats the median random binary partition across the full tested grid, making XOR the paper's clearest instance of a substrate-native logical distinction. NOT, AND, and reversible XOR embeddings exhibit distinct closure regimes rather than a single uniform notion of Boolean success. When a reversible CNOT embedding is viewed through an erased XOR output, closure degrades sharply and unretained input information increases by more than an order-one factor, showing that reversible embedding and erased output are not the same logical object. Finally, a minimal discovery pipeline recovers stable packaged bits and reconstructs a NOT gate directly from transition structure, showing that the framework can find logic rather than merely redescribe it.

\paragraph{Contributions.}
This paper makes three contributions. First, it gives a layer-relative definition of logic in the Six Birds framework: a logic layer is an audited closure of packaged predicates under induced dynamics. Second, it introduces operational diagnostics that distinguish genuine logic layers from projected symbolic descriptions. Third, it demonstrates the framework in controlled finite laboratories, yielding a parity-robustness result, a reversible-versus-erased gate separation, and an unsupervised recovery result for both packaged bits and a NOT gate.

The paper is not a general survey of physical computation, nor a claim that Boolean logic is the only endpoint of coarse-graining. On the contrary, the viewpoint developed here is pluralistic: the logic a layer ``has'' depends on which closures it can stably support. In particular, complement, negation, and even the availability of binary truth carriers are empirical achievements rather than automatic givens. What we do show is that once closure, packaging, and audit are treated as first-class objects, logic becomes something a substrate can earn, lose, or support only partially.

The remainder of the paper develops this claim in stages. Section~\ref{sec:framework} formalizes the notion of a logic layer through packaging, definability, and induced closure. Section~\ref{sec:diagnostics-methods} introduces the diagnostics and the controlled finite laboratories used throughout. Section~\ref{sec:results-parity} shows that parity is an unusually robust coarse variable. Section~\ref{sec:results-phase} studies gate closure regimes for NOT, CNOT, and AND across noise, staging, and timescale stress. Section~\ref{sec:results-reversible} compares reversible embedding to erased XOR and quantifies the resulting accounting differences. Section~\ref{sec:results-discovery} shows that bits and a NOT gate can be recovered directly from transition structure. Section~\ref{sec:discussion} closes by returning to the broader question of what it means for a layer to have a logic at all.
